\documentclass[conference]{IEEEtran}
\IEEEoverridecommandlockouts

\usepackage{cite}
\usepackage{amsmath,amssymb,amsfonts}
\usepackage{graphicx}
\usepackage{textcomp}
\usepackage{xcolor}
\usepackage{booktabs}
\usepackage{array}
\usepackage{url}
\usepackage{float}
\usepackage{placeins}
\newcommand{\tighttables}{%
  \setlength{\textfloatsep}{9pt plus 2pt minus 2pt}%
  \setlength{\floatsep}{8pt plus 2pt minus 2pt}%
  \setlength{\intextsep}{8pt plus 2pt minus 2pt}%
}

\begin{document}
\tighttables

% ---------------------------------------------------------------
% TITLE AND AUTHORS
% ---------------------------------------------------------------
\title{To Anchor or to Deliberate? AI Advice Timing in
       Cost-Sensitive Decision Support}

\author{%
  \IEEEauthorblockN{Anna Asatryan}
  \IEEEauthorblockA{%
    \textit{Akian College of Science and Engineering}\\
    \textit{American University of Armenia}\\
    Yerevan, Armenia\\
    aasatryan04@gmail.com}
  \and
  \IEEEauthorblockN{Karen Hovhannisyan (Supervisor)}
  \IEEEauthorblockA{%
    \textit{Akian College of Science and Engineering}\\
    \textit{American University of Armenia}\\
    Yerevan, Armenia\\
    khovhannisyan@aua.am}
}

\maketitle

% ---------------------------------------------------------------
% ABSTRACT  
% ---------------------------------------------------------------

\begin{abstract}
Predictive models used in decision support are often evaluated as
standalone classifiers, but deployment outcomes depend on how people use
model outputs under task constraints. This paper tests whether AI-advice
timing changes human--AI decision quality in a cost-sensitive loan task.
A web-based within-subject experiment compared three protocols---no~AI,
AI-first, and human-first---using a calibrated logistic regression model
held fixed across conditions. One hundred participants completed 1,800
scored trials under asymmetric penalties: approving a defaulter cost five
units and rejecting a creditworthy applicant cost one unit. Both
AI-assisted protocols significantly reduced mean trial cost relative to
no~AI: AI-first by 0.210 units ($p = 0.017$, $d = 0.243$) and
human-first by 0.298 units ($p < 0.001$, $d = 0.385$). The direct timing
contrast was not statistically decisive ($\Delta = 0.088$, $p = 0.178$,
$d = 0.136$). In human-first trials, AI exposure corrected 90 initial
errors versus 32 it introduced. Weight-of-advice analysis showed
polarized reliance: 52.5\% of trials showed no probability adjustment,
while adjusted trials moved strongly toward the AI estimate (median WOA
$= 0.869$). Although the sample was predominantly young non-experts in
an 18-trial session, limiting external validity to professional settings,
the findings support a system-level view: realized decision value depends
not only on model quality but also on interaction structure, cost
framing, and selective revision behavior.
\end{abstract}

\begin{IEEEkeywords}
human--AI interaction, decision support systems, advice timing,
reliance, cost-sensitive evaluation, behavioral experiment
\end{IEEEkeywords}

% ---------------------------------------------------------------
\section{Introduction}
\label{sec:intro}
% ---------------------------------------------------------------

AI decision-support systems are often justified using offline model
metrics such as accuracy, AUC, or calibration error. These measures are
necessary but not sufficient for evaluating deployed systems. In many
decision contexts the model does not act autonomously; it advises a
human decision-maker. The final decision is therefore produced by an
interaction between the model output, the user, the interface, and the
task incentives. A model can be statistically valid yet operationally
ineffective if users ignore correct advice, over-follow incorrect advice,
or apply the advice using the wrong decision threshold.

This paper studies one interaction-design choice that is often treated
as a single unexamined default: \emph{when} should AI advice be shown?
In an AI-first workflow, the model output appears before the user forms
a judgment. This may improve decision quality by supplying calibrated
information early, but it also risks anchoring the user on the model
output and reducing independent evaluation~\cite{rastogi2022deciding}.
In a human-first workflow, the user first forms an initial judgment and
only then sees the AI prediction. This structure can function as a
cognitive forcing mechanism~\cite{bucinca2021trust}, preserving
independent reasoning before advice exposure, but it may also create
inertia that makes users less willing to adjust when the model is
correct.

The question matters because prior work has studied algorithmic advice,
overreliance, and cognitive forcing, but has less often isolated advice
timing itself. A probability shown before judgment becomes an initial
reference point; the same probability shown after judgment becomes advice
that can be accepted, ignored, or used to revise. Thus, AI-first and
human-first are distinct interaction structures even when the model and
displayed information are held constant.

The study addresses three research questions. RQ1 asks whether AI
probability support improves cost-sensitive decision quality relative to unaided human judgment. RQ2 asks whether advice timing changes realized decision quality when the model, cases, and information content are held constant. RQ3 asks whether, in the human-first workflow, AI exposure corrects more initial errors than it introduces and how strongly users move toward the AI estimate. 

We evaluate advice timing in a loan approval task that exemplifies the
human-centered AI research agenda~\cite{steyvers2024three}.
Participants reviewed loan application profiles and made binary
approve/reject decisions under asymmetric costs reflecting the fact that
approving a defaulting loan is substantially more costly than rejecting a
creditworthy borrower~\cite{thomas2009consumer}. The AI system displayed
a predicted probability of default derived from a calibrated logistic
regression model. The domain provides a structured environment for
studying human--AI decision support; no claims are made about
professional lending practice.

The study makes three contributions. First, it provides empirical
evidence from a completed within-subject experiment comparing no~AI,
AI-first, and human-first protocols while holding the predictive model
and case set constant across conditions. Second, it evaluates decisions
using a cost-sensitive benchmark rather than raw accuracy, exposing
cost-relevant behavioral effects that accuracy metrics conceal. Third,
the human-first protocol enables direct measurement of revision behavior,
including whether AI exposure corrects initial errors or introduces new
ones.

The main finding is that AI support improved decisions relative to no~AI
under both cost and accuracy metrics with medium effect sizes. Human-first
support produced the lowest cost descriptively and showed clear
correction behavior, but the direct timing difference versus AI-first
was not statistically decisive at the achieved sample size.

% ---------------------------------------------------------------
\section{Literature Review}
\label{sec:lit}
% ---------------------------------------------------------------

\subsection{The Model--System Performance Gap}

A foundational challenge motivating this work is that deploying an
accurate predictive model does not automatically improve decisions.
Kleinberg et al.~\cite{kleinberg2018human} demonstrated this through a
retrospective policy simulation of bail decisions, showing via
counterfactual analysis that an algorithm substantially outperformed
judges at predicting flight risk. Their analysis was conducted on
historical data rather than in a live human-AI interaction setting,
leaving open whether such gains would materialize when decision-makers
actually use algorithmic predictions.

Green and Chen~\cite{green2019principles} addressed this question
directly through a controlled behavioral experiment. They found that
algorithmic risk scores reliably improved participant accuracy over
unaided baselines in most conditions, but participants failed to match
the model's standalone accuracy, could not evaluate prediction quality,
and did not calibrate their reliance based on the model's actual
performance. We adopt their algorithm-in-the-loop framework as our
central lens, treating the gap between model accuracy and realized
decision quality as the primary quantity of interest.

Mozannar and Sontag~\cite{mozannar2020consistent} approached this gap
from a machine learning perspective, formalizing the learning-to-defer
problem: given that a human will make the final decision, how should the
model be trained to allocate decision authority optimally? Their framework
treats human performance as a function of case difficulty and formalizes
the system-level objective that this paper studies empirically. We
complement their theoretical contribution with behavioral data on how
interaction structure, rather than model architecture, determines whether
the potential for complementary performance is realized.

\subsection{Reliance, Overreliance, and Cognitive Forcing}

The human factors literature documents the tendency toward automation
bias: the uncritical acceptance of automated recommendations even when
human judgment would be superior~\cite{parasuraman2010complacency}. In
AI-assisted decision settings this manifests as overreliance---following
the model when it is incorrect---which is systematically documented
across clinical, judicial, and financial assessment
domains~\cite{bucinca2021trust,schemmer2023appropriate}.

Bu\c{c}inca et al.~\cite{bucinca2021trust} showed that cognitive forcing
functions---design interventions requiring users to form an independent
judgment before receiving AI advice---can substantially reduce
overreliance. Their finding is directly relevant to our human-first
protocol, which embeds a structural cognitive forcing mechanism by
collecting an initial decision before revealing the AI prediction.
Their work, however, does not systematically vary the timing of AI advice
across a common case set while controlling for model quality and case
difficulty, which is the precise contribution of this study.

Schemmer et al.~\cite{schemmer2023appropriate} proposed the
Appropriateness of Reliance (AoR) framework, decomposing reliance into
two dimensions: correct AI reliance (following the AI when it is correct)
and correct self-reliance (overriding the AI when it is wrong). This
two-dimensional view is essential: a participant who always follows the AI
scores high on the first dimension but near zero on the second. We
operationalize a related decomposition in our human-first correction
analysis, distinguishing beneficial reliance, harmful override,
over-reliance, and beneficial override across the full trial set.

\subsection{Timing, Anchoring, and Advice Taking}

Advice timing is connected to anchoring. If the AI output is shown first,
it may become the reference point around which the human forms a
judgment. Rastogi et
al.~\cite{rastogi2022deciding} showed mathematically that anchoring on AI
predictions can degrade collaborative performance and proposed enforced
time delays as a de-anchoring mechanism. Their framework informs our
measurement approach, though we manipulate the ordering of information
rather than enforcing delays.

Steyvers and Kumar~\cite{steyvers2024three} identify sequential
interaction and advice timing as important unresolved challenges in
AI-assisted decision making, noting that the time-delayed paradigm has
shown promise for reducing anchoring but that ordering effects remain
understudied in fully within-subject designs. Bansal et
al.~\cite{bansal2021does} showed that complementary human--AI team
performance depends on task structure and the conditions under which
humans choose to override the model, suggesting that interaction design
mediates whether potential complementarity is realized.

The weight of advice (WOA) metric~\cite{bonaccio2006advice,harvey1997taking}
measures the fraction of movement from an initial estimate toward an
advisor's estimate and is only observable when an initial judgment is
collected before advice exposure. This makes WOA especially suited to the
human-first condition. Prior WOA studies typically find conservative
advice-taking with participants adjusting less than the normatively optimal
amount, but the direction and magnitude of adjustment can vary
substantially with advisor type and task domain.

\subsection{Confidence, Calibration, and Uncertainty Communication}

How AI uncertainty is communicated affects both reliance behavior and
participant calibration. Zhang et al.~\cite{zhang2020effect} demonstrated
that displaying model confidence alongside predictions can improve trust
calibration---participants become more discriminating about when to follow
AI advice---but that this benefit depends on the confidence signal being
genuinely informative. When confidence is poorly calibrated, displaying it
can increase overreliance.

This motivates our decision to use a probability-based AI representation
derived from isotonic calibration~\cite{niculescu2005predicting} rather
than a binary recommendation. The interface displayed the model's
predicted default probability as the AI output, while the task's cost rule
was introduced during instructions and remained available throughout the
experiment. This allowed us to study probability anchoring and revision
behavior without presenting an explicit approve/reject recommendation. The
model's mean confidence (probability assigned to the predicted class)
averages 0.817 when correct and 0.731 when incorrect, yielding a
discriminative gap that is large enough to be informative but small
enough to create a genuinely non-degenerate reliance environment.

\subsection{Case Difficulty and Complementary Performance}

The relationship between task difficulty and human-AI team performance
is central to achieving complementary outcomes. Bansal et
al.~\cite{bansal2021does} found that while AI augmentation produced
complementary improvements over unaided human performance, AI
explanations did not increase these gains and sometimes increased
inappropriate reliance. A key implication is that complementary
performance requires the human to selectively override incorrect AI
predictions, which demands that human and AI errors be non-overlapping.

These findings motivate an experimental design in which case difficulty
and model correctness are both represented. Including cases where the
model is correct and incorrect creates opportunities to observe both
beneficial reliance and beneficial override, while medium-difficulty
cases are especially important because model--optimal disagreement is
most consequential near the cost-sensitive decision threshold.

Credit decisions provide a well-structured domain for cost-sensitive
human-AI research~\cite{thomas2009consumer}. The cost asymmetry between
false negatives (approving defaulters) and false positives (rejecting
creditworthy borrowers) creates a normative framework that is absent in
accuracy-only evaluation~\cite{elkan2001foundations}, enabling direct
comparison between human behavior and a cost-minimizing benchmark.

% ---------------------------------------------------------------
\section{Methodology}
\label{sec:method}
% ---------------------------------------------------------------

\subsection{Prediction Task and Data}

The loan default prediction task was built on LendingClub consumer loan
data~\cite{lendingclub2018}. The prediction target was whether a loan defaulted. Only variables
observable at loan origination were used; post-outcome fields,
payment-history variables, and internal credit-grade labels were excluded
to prevent data leakage. The cleaned modeling dataset contained
1,298,960 loans issued from 2007 to 2018 with a global default rate of
20.1\%.

The data were split temporally. Loans issued before 2016 were used for
model development and calibration (default rate 18.5\%); loans issued
from 2016 onward formed the held-out test period from which experimental
cases were drawn ($n = 479{,}744$; default rate 22.9\%). This split
preserves realistic temporal distribution shift and avoids using future
lending outcomes to build the AI model shown to participants.

The final feature set contains nine variables satisfying four criteria:
available at decision time, understandable to participants without lending expertise, showing a monotonic relationship with default risk in training data, and non-redundant. The retained features are loan amount, term, interest rate, annual income (log-transformed), debt-to-income ratio, revolving
utilization, home ownership, loan purpose, and credit history length.
Redundant variables such as installment (functionally determined by loan
amount, term, and rate) and low-signal variables such as delinquency
count were excluded. Numeric features use median imputation and
z-score standardization; categorical features use mode imputation and
one-hot encoding.

\subsection{AI Model}

Three candidate classifiers were evaluated on a shared pipeline: logistic
regression ($L_2$ regularization, $C = 0.1$), XGBoost (300 estimators,
depth 4), and histogram gradient boosting (300 iterations, depth 4).
Hyperparameters were fixed rather than grid-searched, as the objective is
a stable, defensible probabilistic advisor rather than maximum leaderboard
performance.

A ten-run bootstrap stability analysis on the training set showed Brier
score standard deviations below 0.001 for all three model classes,
confirming that the model comparison was stable rather than a
single-split artifact. On the held-out post-2016 test set, Brier score
differences between logistic regression and the boosting alternatives were
at most 0.0012 (Table~\ref{tab:model}). Logistic regression with isotonic
calibration was selected as the decision-support model for its transparent
coefficient structure, stable training behavior, and negligible
performance difference from more complex alternatives.

Calibration was applied after fitting using isotonic regression on a
held-out calibration split (126,135 pre-2016 loans), following the
prefit calibration procedure to avoid refitting the base
model~\cite{niculescu2005predicting}. On the post-2016 test set,
isotonic calibration produces ECE = 0.033. Calibration quality varies
by difficulty tier (easy: 0.025; medium: 0.040; hard: 0.050), reflecting
higher uncertainty in the ambiguous mid-range where the model has less
predictive leverage.

\begin{table}[H]
  \centering
  \caption{Held-Out Model Performance on the Post-2016 Test Set ($n = 479{,}744$)}
  \label{tab:model}
  \scriptsize
  \setlength{\tabcolsep}{3.5pt}
  \begin{tabular}{lccc}
    \toprule
    \textbf{Model} & \textbf{AUC} & \textbf{Brier} & \textbf{ECE} \\
    \midrule
    Logistic regression   & 0.691 & 0.164 & 0.033 \\
    XGBoost               & 0.697 & 0.163 & 0.034 \\
    Hist. grad. boosting  & 0.697 & 0.163 & 0.033 \\
    \bottomrule
  \end{tabular}

  \vspace{2pt}
  \parbox{\columnwidth}{\scriptsize
    Bootstrap 95\% CI for Logistic AUC: [0.690, 0.693];
    Brier: [0.163, 0.165].
  }
\end{table}

\subsection{Case Selection and Difficulty Stratification}

Eighteen scored cases and two practice cases were used. Case difficulty
was defined model-independently using empirical default rates estimated
on training data, avoiding circularity between difficulty and model
uncertainty. For each combination of interest-rate, DTI, and income bin
(3D grid), a Bayesian-smoothed default rate was computed, shrinking
sparse-bin estimates toward the training mean:
\begin{equation}
  \hat{p}^{\,\text{smooth}}_b =
    \frac{n_b \,\bar{p}_b + \alpha\,\bar{p}_{\text{train}}}
         {n_b + \alpha},\quad \alpha = 50.
  \label{eq:smooth}
\end{equation}
Difficulty was then defined as closeness to maximum uncertainty:
\begin{equation}
  \text{difficulty}_i = 1 - \frac{|\hat{p}_i - 0.5|}{0.5},
  \label{eq:difficulty}
\end{equation}
yielding a score in $[0,1]$ partitioned into three tiers (easy, medium,
hard) at the 33rd and 66th percentiles. Smoothing sensitivity analysis
confirmed that tier assignments were stable for $\alpha \in [10, 100]$.

Scored cases were sampled from the post-2016 test set to span 3
difficulty tiers $\times$ 2 model-correctness states = 6 cells, with
3~cases per cell, for a total of 18. Within each cell, cases were
selected to maximize diversity of predicted probability (minimum pairwise
gap of 0.02) and loan purpose. The selected cases span predicted
probabilities from 0.055 to 0.946 (mean 0.441). All case files were
locked before participant recruitment and have not been modified since.

\subsection{Interaction Protocols}
\label{sec:protocols}

The experiment used three within-subject protocols. In the \emph{no-AI}
protocol, only the loan profile was displayed; participants entered a
final binary decision and a probability estimate. In the
\emph{AI-first} protocol, the model's predicted default probability was
displayed before the participant's decision. In the \emph{human-first}
protocol, participants first submitted an initial binary decision and
probability estimate; the AI probability was then revealed, after which
participants had the opportunity to revise both responses before final
submission. The asymmetric cost rule was introduced before scored trials
and remained accessible throughout the experiment.

Table~\ref{tab:protocols} summarizes the three protocol structures.
The human-first protocol is the only condition that enables pre/post
revision measurement and WOA analysis; the initial and final submissions
are stored separately in the database.

\begin{table}[H]
  \centering
  \caption{Interaction Protocols Used in the Within-Subject Experiment}
  \label{tab:protocols}
  \small
  \begin{tabular}{>{\raggedright\arraybackslash}p{0.24\columnwidth}
                  >{\raggedright\arraybackslash}p{0.68\columnwidth}}
    \toprule
    \textbf{Protocol} & \textbf{Participant flow} \\
    \midrule
    No~AI & Loan profile; final probability estimate; approve/reject decision. \\
    AI-first & Loan profile with AI probability; final probability estimate; approve/reject decision. \\
    Human-first & Loan profile; initial estimate and decision; AI probability; final estimate and decision. \\
    \bottomrule
  \end{tabular}
\end{table}

\subsection{Experiment Platform and Reproducibility}

The experiment was implemented as a web application in Streamlit with a
Supabase backend. The platform enforced protocol-specific trial flow,
stored participant records, trial-level decisions, probability estimates,
protocol labels, case identifiers, and timing metadata at the individual
trial level. This implementation allows behavioral data to be reconstructed and linked to frozen experimental cases using anonymous participant identifiers. The study did not collect direct personally identifying information; exported
data contain participant-level responses, protocol assignments, timing
metadata, and outcomes needed for reproducible analysis.

Reproducibility was handled through locked artifacts. The case-design
notebook produced candidate pools and working outputs to a build
directory; the official experimental snapshot was stored separately as
frozen case files with a lock file (\texttt{cases.lock.json}) and
integrity validation. The submitted repository separates the official
behavioral reproduction from the optional upstream raw-data rebuild. The command \texttt{python run.py} regenerates the final reported tables, figures, and summary from
\texttt{data/experiment\_exports/} and \texttt{artifacts/frozen/}. All paper figures are generated programmatically by the reproduction
pipeline and copied into the paper source without manual editing. The raw file
\texttt{data/raw/loan.csv} is needed only for the optional provenance
command \texttt{python run.py --mode rebuild}, which writes to
\texttt{artifacts/build/} and does not modify the frozen experimental
boundary.

\subsection{Participants}

A total of 158 participants enrolled; 100 completed all scored trials
(63.3\% completion rate). A chi-square balance test confirmed that the
three counterbalancing groups were not meaningfully unbalanced among
completers ($\chi^2 = 0.560$, $p = 0.756$): 30 in group~1, 36 in
group~2, and 34 in group~3. Each completed participant contributed
18~scored trials, yielding 1,800 scored decisions in total.

The sample was predominantly young (78\% aged 18--24; 14\% aged 25--34)
and student-based (56\% bachelor's in progress; 28\% master's or higher).
No participant reported professional lending experience. Participants
completed a three-question comprehension quiz before scored trials;
80.2\% passed on the first attempt.

\subsection{Outcome Measures and Statistical Analysis}

The primary outcome was mean trial cost under the asymmetric cost function:
\begin{equation}
  \text{cost}(d,\,y) = \begin{cases}
    C_{FN} = 5 & \text{if } d = \text{approve},\ y = \text{default} \\
    C_{FP} = 1 & \text{if } d = \text{reject},\ y = \text{paid}     \\
    0           & \text{otherwise,}
  \end{cases}
  \label{eq:cost}
\end{equation}
where $d$ is the participant's binary decision and $y$ is the true loan
outcome. This ratio implies the normative cost-minimizing threshold
$\tau = C_{FP} / (C_{FP} + C_{FN}) = 1/6 \approx 0.167$~\cite{elkan2001foundations}:
approve if and only if predicted default probability is below $\tau$.
Accuracy was a secondary outcome.

Weight of advice was computed for human-first trials as:
\begin{equation}
  \text{WOA}_i =
    \frac{p^{\text{final}}_i - p^{\text{initial}}_i}
         {p^{\text{AI}}_i - p^{\text{initial}}_i}.
  \label{eq:woa}
\end{equation}
Trials where $|p^{\text{AI}}_i - p^{\text{initial}}_i| < 0.01$ were
excluded (21 of 600; 3.5\%) as adjustment toward the AI estimate was
undefined in such cases. Values were clipped to $[-1, 2]$ for stability.

The primary inference used paired $t$-tests on per-participant protocol
means, which directly follows from the within-subject design and avoids
pseudo-replication from treating 1,800 trials as independent. Three
planned contrasts were evaluated for each outcome: no~AI vs.\ AI-first,
no~AI vs.\ human-first, and AI-first vs.\ human-first. A trial-level
linear mixed model was attempted as a robustness check but produced a
singular covariance matrix, indicating the random-effects structure was
over-parameterized for the within-participant variance structure; paired
$t$-tests are therefore the sole primary inference. Effect sizes are
reported as Cohen's $d$.

% ---------------------------------------------------------------
\section{Results}
\label{sec:results}
% ---------------------------------------------------------------

\subsection{Decision Cost}

Table~\ref{tab:cost} shows the primary result. Mean trial cost was
1.222 in the no-AI condition, 1.012 with AI-first, and 0.923 with
human-first. Both AI-assisted protocols significantly reduced cost
relative to no~AI. Human-first produced the lowest cost numerically.
The direct timing contrast between AI-first and human-first was not
statistically decisive: $\Delta = 0.088$ (95\% CI: $[-0.041,\, 0.218]$;
$t = 1.356$, $p = 0.178$, $d = 0.136$).

\begin{table}[H]
  \centering
  \caption{Primary Outcome: Mean Trial Cost by Protocol ($N=100$)}
  \label{tab:cost}
  \footnotesize
  \setlength{\tabcolsep}{3.2pt}
  \begin{tabular}{@{}lcccc@{}}
    \toprule
    \textbf{Comparison} & \textbf{A} & \textbf{B}
      & $\boldsymbol{\Delta}$ & $\boldsymbol{p}$ \\
    \midrule
    No~AI vs.\ AF & 1.222 & 1.012 & $+0.210$ & 0.017 \\
    No~AI vs.\ HF & 1.222 & 0.923 & $+0.298$ & $<0.001$ \\
    AF vs.\ HF    & 1.012 & 0.923 & $+0.088$ & 0.178 \\
    \midrule
    Model-optimal ($\tau=1/6$)$^\dagger$ & --- & 0.333 & --- & --- \\
    \bottomrule
  \end{tabular}

  \vspace{2pt}
  \parbox{\columnwidth}{\footnotesize
    A and B are protocol means; $\Delta=A-B$, so positive values mean A is worse.
    AF = AI-first; HF = human-first.
    $^\dagger$Observed realized cost from applying $\tau=1/6$ to AI probabilities.
    95\% CIs: no~AI vs.\ AF (0.038, 0.382); no~AI vs.\ HF (0.144, 0.452).
  }
\end{table}

\subsection{Decision Accuracy}

Table~\ref{tab:accuracy} shows accuracy contrasts. Mean accuracy was
55.2\% without AI, 63.5\% with AI-first, and 65.0\% with human-first.
Both AI-assisted protocols significantly improved accuracy relative to
no~AI; the AI-first vs.\ human-first contrast was not significant.

\begin{table}[H]
  \centering
  \caption{Secondary Outcome: Accuracy by Protocol ($N=100$)}
  \label{tab:accuracy}
  \begin{tabular}{lcccc}
    \toprule
    \textbf{Comparison} & \textbf{Mean A} & \textbf{Mean B}
      & $\boldsymbol{\Delta}$ & $\boldsymbol{p}$ \\
    \midrule
    No~AI vs.\ AI-first        & 0.552 & 0.635 & $-0.083$ & $<0.001$    \\
    No~AI vs.\ Human-first     & 0.552 & 0.650 & $-0.098$ & $<0.001$    \\
    AI-first vs.\ Human-first  & 0.635 & 0.650 & $-0.015$ & $0.408$     \\
    \bottomrule
    \multicolumn{5}{l}{\footnotesize $\Delta$ = Mean~A $-$ Mean~B; negative = A is less accurate.}\\
    \multicolumn{5}{l}{\footnotesize Effect sizes: no~AI vs.\ AF $d=-0.392$; no~AI vs.\ HF $d=-0.467$.}
  \end{tabular}
\end{table}

The AI benefit by difficulty tier was concentrated in easy cases
(Table~\ref{tab:diff}). Because difficulty was defined using
training-data empirical default uncertainty rather than participant-rated
difficulty, these tiers should be interpreted as statistical case strata
rather than subjective difficulty ratings.

\begin{table}[H]
  \centering
  \caption{Accuracy by Protocol and Difficulty Tier}
  \label{tab:diff}
  \scriptsize
  \setlength{\tabcolsep}{3.2pt}
  \begin{tabular}{@{}lccc@{}}
    \toprule
    \textbf{Tier} & \textbf{No~AI} & \textbf{AI-first}
      & \textbf{Human-first} \\
    \midrule
    Easy   & 59.0\% & 81.0\% & 82.0\% \\
    Medium & 58.0\% & 58.5\% & 62.5\% \\
    Hard   & 48.5\% & 51.0\% & 50.5\% \\
    \bottomrule
  \end{tabular}

  \vspace{1pt}
  \parbox{0.82\columnwidth}{\scriptsize
    Each cell contains 200 trials.
  }
\end{table}

The protocol-level cost and accuracy pattern is summarized in
Fig.~\ref{fig:cost_accuracy}.

\begin{figure}[H]
  \centering
  \includegraphics[width=0.96\columnwidth]{figures/cost_accuracy_by_protocol.png}
  \caption{Mean trial cost (left) and accuracy (right) by protocol across
    100 participants. Lower cost and higher accuracy are better. Confidence
    intervals for all pairwise contrasts are reported in
    Tables~\ref{tab:cost} and~\ref{tab:accuracy}.}
  \label{fig:cost_accuracy}
\end{figure}

\subsection{Human-First Correction Behavior}

Across 600 human-first trials, 300 (50.0\%) stayed correct after AI
exposure, 90 (15.0\%) improved from initially wrong to finally correct,
178 (29.7\%) stayed wrong, and 32 (5.3\%) worsened from initially
correct to finally wrong (Table~\ref{tab:correction}). The net gain was
$+58$ trials (9.7 percentage points). Among the 122 trials where
correctness changed, a one-sided exact binomial test confirmed that
improvements significantly outnumbered worsenings ($90$ vs.\ $32$;
$p < 0.0001$).

\begin{table}[H]
  \centering
  \caption{Human-First Pre/Post Correction Matrix ($n = 600$ trials)}
  \label{tab:correction}
  \begin{tabular}{lcc}
    \toprule
                    & \textbf{Final correct} & \textbf{Final wrong}  \\
    \midrule
    \textbf{Initial correct} & 300~~(50.0\%)  & ~~32~~(5.3\%)   \\
    \textbf{Initial wrong}   & ~~90~~(15.0\%) & 178~~(29.7\%)   \\
    \bottomrule
  \end{tabular}
\end{table}

The same pre/post correction structure is visualized in
Fig.~\ref{fig:correction_matrix}.

\begin{figure}[H]
  \centering
  \includegraphics[width=0.76\columnwidth]{figures/human_first_correction_matrix.png}
  \caption{Human-first correction matrix. Green corresponds to trials that
    improved after AI exposure, red to trials that worsened, and gray to
    trials whose correctness status did not change.}
  \label{fig:correction_matrix}
\end{figure}

\subsection{Weight of Advice and Reliance Heterogeneity}

Of 579 WOA-eligible human-first trials, 304 (52.5\%) showed essentially
no probability movement (absolute change $< 0.01$). Among the 275
adjusted trials, median WOA was 0.869 and mean WOA was 0.861
(Fig.~\ref{fig:woa}). The overall WOA distribution (mean 0.409; median
0.000) conceals a bimodal pattern: non-adjustment versus near-complete
adoption of the AI estimate among those who did adjust.

\begin{figure}[H]
  \centering
  \includegraphics[width=0.95\columnwidth]{figures/woa_distribution.png}
  \caption{(a) Adjustment incidence among 579 WOA-eligible human-first
    trials. (b) WOA distribution among the 275 adjusted trials
    (those with $|p^{\text{final}} - p^{\text{initial}}| \geq 0.01$).
    Dashed line: median WOA $= 0.869$; dotted line: WOA $= 1$ (full
    adoption of AI estimate).}
  \label{fig:woa}
\end{figure}

Individual-level AI benefit on accuracy (mean $= 0.091$, SD $= 0.191$)
further confirms substantial between-participant heterogeneity. The
distribution has a negative tail: approximately one-third of participants
showed lower accuracy under AI-supported protocols than under no~AI,
suggesting a subgroup for whom AI support was not beneficial.

\FloatBarrier

% ---------------------------------------------------------------
\section{Analysis and Validation}
\label{sec:analysis}
% ---------------------------------------------------------------

\subsection{Did the Study Objectives Succeed?}

The study had three objectives: (1) quantify the effect of AI support on
cost-sensitive decisions; (2) test whether timing of AI advice modulates
decision quality; (3) characterize reliance behavior in the human-first
condition. On objective~1, both AI protocols reduced cost and improved
accuracy with medium effect sizes, and the human-first correction analysis
confirmed net-positive AI-induced corrections ($p < 0.0001$, binomial).
On objective~2, the timing contrast was directional but not statistically
decisive at $N = 100$ ($p = 0.178$; 95\% CI includes zero). On
objective~3, WOA analysis documented strongly bimodal reliance, and the
correction matrix decomposed AI-induced changes into beneficial and
harmful revisions.

\subsection{Alignment with Prior Work}

The magnitude of AI benefit aligns with prior algorithm-in-the-loop
findings that decision support can improve unaided human judgment but
rarely realizes the full standalone model potential~\cite{green2019principles,bansal2021does}.
The concentration of benefit in easy cases mirrors Bansal et al.'s
observation~\cite{bansal2021does} that complementary human--AI performance
depends on task structure: where the AI is confident and correct, the
benefit is large; where the AI is ambiguous, humans provide limited
additional signal.

The human-first correction behavior is consistent with the cognitive
forcing function literature~\cite{bucinca2021trust}. Requiring initial
judgment before AI exposure produced measurably more corrections than
degradations ($90:32$). However, 178 trials (29.7\%) remained wrong even
after AI exposure, indicating that the forcing function alone is
insufficient for cases where participants lacked the domain knowledge to
integrate the AI signal correctly.

The WOA distribution is consistent with discrete switching behavior
documented in advice-taking research~\cite{bonaccio2006advice,harvey1997taking}. The bimodal pattern---ignoring versus fully
adopting---is less consistent with the conservative-updating framework,
which predicts intermediate WOA values. This suggests that participants
in this domain treated AI advice as a categorical override signal rather
than a Bayesian update.

\subsection{Comparison with Benchmark Strategies}

Table~\ref{tab:bench} compares observed human protocol means with
reference decision strategies on the 18-case set. Human performance under
all protocols substantially exceeded the always-approve and random
baselines, indicating that participants were genuinely discriminating.
However, all human conditions fell substantially above the model-optimal
benchmark (0.333), a gap that even human-first support ($0.923$) did not
close. This gap reflects both imperfect calibration and sub-optimal
threshold application, not model quality.

\begin{table}[H]
  \centering
  \caption{Cost Benchmarks on the 18-Case Set}
  \label{tab:bench}
  \begin{tabular}{lc}
    \toprule
    \textbf{Strategy}                      & \textbf{Cost} \\
    \midrule
    Always approve                          & 2.500 \\
    Random (50\%)                           & 1.500 \\
    Model at $p < 0.5$                      & 1.833 \\
    No~AI (observed)                        & 1.222 \\
    AI-first (observed)                     & 1.012 \\
    Human-first (observed)                  & 0.923 \\
    Always reject                           & 0.500 \\
    Model-optimal ($\tau = 1/6$, observed)  & 0.333 \\
    \bottomrule
  \end{tabular}
\end{table}

\subsection{Robustness Checks}

\textbf{Group-3 sensitivity.} Excluding the 34 participants in
counterbalancing group~3 (reduced sample $N = 66$) preserved all
substantive conclusions. The no~AI vs.\ AI-first contrast remained
significant ($t = -4.959$, $p < 0.0001$, $d = -0.610$) and no~AI vs.\
human-first remained significant ($p < 0.0001$, $d = -0.660$). The
AI-first vs.\ human-first contrast remained non-significant ($p =
0.658$).

\textbf{Non-parametric check.} Wilcoxon signed-rank tests on
per-participant accuracy means produced the same pattern of significant
and non-significant contrasts as the paired $t$-tests, indicating that
the results are not driven by distributional assumptions.

\textbf{Model selection stability.} Bootstrap resampling over 10 runs
confirmed Brier score standard deviations below 0.001 across all three
candidate models. The final logistic regression model was selected before
any participant data were collected; model evaluation on the post-2016
test set served as a final reporting diagnostic only.

\subsection{Normative Deviation and Threshold Analysis}

Normative deviation from the cost-sensitive threshold was 26.0\% for
easy cases, 73.3\% for medium cases, and 4.7\% for hard cases. The
medium tier is the region where the optimal rule ($\tau = 1/6$) requires
rejection but a naive accuracy-maximizing rule at 0.5 would approve.
Participants deviated most frequently in this region, consistent with
application of a 0.5 threshold rather than the cost-calibrated $\tau$.


Implicit decision threshold estimates, derived by fitting logistic
regression to no-AI trial decisions, yielded valid estimates for 96 of
100 participants (mean $= 0.355$, median $= 0.299$), substantially above
the normative $\tau = 0.167$. This confirms that participants
systematically approved loans at predicted default probabilities that
exceed the cost-minimizing threshold, contributing to the persistent gap
from the model-optimal benchmark.

Decision consistency between final binary decisions and probability
estimates was 84.6\% at $\tau = 0.167$ and 79.1\% at $\tau = 0.5$,
indicating that participants' own probability estimates were more
consistent with the cost-sensitive rule than their actual decisions.
The case-level cost pattern is shown in Fig.~\ref{fig:case_scatter}.

\begin{figure}[htbp]
  \centering
  \includegraphics[width=0.92\columnwidth]{figures/case_risk_cost_scatter.png}
  \caption{Case-level mean cost by AI-predicted default probability. Each
    point is one of the 18 scored cases aggregated across participants;
    the dashed vertical line marks the cost-sensitive threshold
    $\tau=1/6$.}
  \label{fig:case_scatter}
\end{figure}

\subsection{Validation Limits}

The validation evidence has three main limits: the sample was young and
non-expert, each participant completed only 18 scored trials, and the
mixed-effects model did not converge. Therefore, the results should be
read as controlled evidence about interaction design, with the fuller
generalizability and power limits discussed in Section~\ref{sec:discussion}.

% ---------------------------------------------------------------
\section{Discussion and Conclusion}
\label{sec:discussion}
% ---------------------------------------------------------------

\subsection{Main Contributions}

This paper makes three empirical contributions to human-centered AI
decision-support research. First, it shows in a completed within-subject
experiment that the same calibrated model improves cost-sensitive
decisions under both AI-first and human-first workflows relative to
unaided judgment. Second, it isolates advice timing while holding the
model, case set, and displayed information constant; human-first produced
the lowest mean cost descriptively, but the direct timing contrast with
AI-first was not statistically decisive. Third, it uses the human-first
condition to measure revision behavior directly, showing that AI exposure
corrected more initial errors than it introduced and that probability
adjustment was highly polarized.

\subsection{Design Implications}

The results support a system-level view of AI decision support. Model
quality alone did not determine realized decision value: even with a
calibrated probability model, human decisions remained far from the
model-optimal benchmark. The gap between no AI (1.222), human-first
(0.923), and model-optimal behavior (0.333) indicates that interface
timing, threshold interpretation, and revision behavior jointly shape how
much predictive value is converted into better decisions.

For applied systems, the results argue against treating advice display as
a neutral interface detail. AI-first workflows are efficient and improved
performance relative to no AI, but they may anchor users before they form
an independent judgment. Human-first workflows preserve independent
deliberation and make revision behavior observable, but they add friction
and can create inertia when users should adjust. A practical system could
therefore reserve human-first interaction for ambiguous, high-stakes, or
low-confidence cases, while using AI-first interaction for clearer cases
where the model signal is strong. This recommendation is directional
rather than prescriptive, because the present sample consisted of
non-expert participants and the timing contrast was not statistically
decisive.

The results also suggest that probability support should be paired with
explicit threshold guidance. Participants often behaved as if the relevant
decision boundary were closer to 0.5, even though the cost-sensitive
threshold was $\tau = 1/6$. Thus, showing calibrated probabilities is not
sufficient if users do not know how those probabilities should translate
into decisions under asymmetric costs. This aligns with prior work on
confidence and explanation, which shows that uncertainty information is
useful only when users can interpret it appropriately~\cite{zhang2020effect}.


\subsection{Limitations}

Several constraints bound interpretation. The sample consisted mainly of
young, non-expert participants (78\% aged 18--24, no professional lending
experience), limiting generalization to loan officers or credit analysts.
The experiment used 18 scored trials per participant, which limited
statistical power for detecting the small AI-first vs.\ human-first
effect ($d = 0.136$); a simulation-based power analysis~\cite{green2016simr}
suggests $N \approx 350$--$450$ would be required for 80\% power at the
observed effect size. The linear mixed model failed to converge, precluding
block-order adjustment; though paired tests are robust, learning effects
within session cannot be formally separated from protocol effects. The AI
advice is a scalar probability without feature attribution; results may
differ under explainable AI paradigms~\cite{zhang2020effect} that provide
feature-level rationales.

Internal validity is strengthened by the within-subject design,
counterbalanced protocol order, locked case set, and fixed AI model, but
it is not complete. Because each participant experienced all three
protocols in a single session, practice and fatigue effects may remain
partially confounded with protocol exposure despite counterbalancing.
Construct validity is also bounded: the task captures structured
probability-assisted decision making, not the full institutional, legal,
and ethical context of real credit allocation. External validity is
therefore limited to similar decision-support settings where users
interpret probabilistic advice under explicit asymmetric costs.

\subsection{Future Directions}

Future work should examine the timing contrast with larger and more
expert samples, more trials per participant, and populations with
professional decision-making experience. The WOA polarization pattern
warrants investigation into individual-level moderators---such as domain
knowledge, numeracy, and prior AI exposure---that predict whether a
participant will adjust or anchor. Hybrid workflow designs that adaptively
apply human-first or AI-first structure based on case difficulty or model
confidence represent a promising applied direction~\cite{mozannar2020consistent}.
Finally, the persistent gap from the model-optimal benchmark across all
human conditions suggests that threshold education and cost-framing
interventions deserve investigation as complements to advice-timing
manipulations.

\subsection{Conclusion}

This paper evaluated AI advice timing in a controlled cost-sensitive
decision-support experiment. One hundred participants completed a
within-subject task comparing no~AI, AI-first, and human-first protocols.
AI support improved decision quality relative to no~AI under both cost
and accuracy metrics. Human-first produced the strongest descriptive
performance and corrected more initial errors than it introduced, but the
direct timing difference between human-first and AI-first was not
statistically decisive. The main conclusion is that human--AI decision
quality is a property of the whole interaction system: a calibrated
model can improve decisions, but the realized gain depends on when advice
is shown, whether users revise appropriately, and whether decisions are
evaluated against the correct cost structure. For human-centered AI
applications requiring high-stakes decision support, model validation
should therefore be paired with behavioral validation of the interaction
workflow through which model outputs reach human judgment.


% ---------------------------------------------------------------
\section*{Acknowledgment}
% ---------------------------------------------------------------

The author thanks Karen Hovhannisyan for supervising this capstone project,
the study participants for their time, and the American University of
Armenia for academic support. The author used Claude AI for language
editing and style refinement of the manuscript. All analyses, results,
interpretations, and final decisions were reviewed and verified by the
author.

% ---------------------------------------------------------------
\appendices
\section{Experiment Platform Implementation}
\label{app:platform}
% ---------------------------------------------------------------

The behavioral experiment was implemented as a Streamlit web application
with Supabase/PostgreSQL persistence. The participant-facing app used a
phase-router structure: the entry point established a browser session
identifier, resumed incomplete sessions when the same identifier was found
in Supabase, routed completed sessions to an already-completed screen, and
selected the appropriate UI screen from a phase-to-screen router. The main
participant flow was consent, demographics, glossary and instructions,
comprehension quiz, two practice trials with feedback, three scored blocks
of six trials each, reflection questions, trust rating, performance
summary, and thank-you screen.

Each scored block used exactly one protocol
(\texttt{no\_ai}, \texttt{ai\_first}, or \texttt{human\_first}), and each
participant saw all three protocols exactly once. Protocol order was
assigned using the locked Latin-square rotation for \texttt{group\_1},
\texttt{group\_2}, and \texttt{group\_3}. The deployed platform loaded the
18 scored cases, two practice cases, and protocol rotation from frozen CSV
files rather than manually embedding experimental stimuli in code.

The Supabase schema used three principal tables:
\texttt{participants}, \texttt{trials}, and
\texttt{quiz\_responses}. The \texttt{participants} table stored group
assignment, completion state, session-resumption fields, total and optimal
cost, trust rating, and post-trial reflection responses. The
\texttt{trials} table stored one row per practice or scored trial,
including protocol, case identifier, AI prediction, true outcome,
initial/final judgments, probability estimates, timing fields, and
correctness indicators. The \texttt{quiz\_responses} table logged
comprehension-quiz answers by participant, attempt, and question. Trial
rows used a uniqueness constraint on \texttt{(participant\_id,
trial\_index)} to prevent duplicate submissions.

The implementation separated responsibilities across modules: the app
entry point handled session bootstrap and routing, a database module
isolated Supabase reads/writes, and an experiment-logic module handled
group assignment, deterministic trial sequencing, and scoring. A demo mode
was available for presentation without storing study responses. After data
collection, completed participant, trial, and quiz tables were exported
from Supabase as CSV files under \texttt{data/experiment\_exports/}, which
are the direct inputs to the official reproduction pipeline.

% ---------------------------------------------------------------
\bibliographystyle{IEEEtran}
\bibliography{references}

\end{document}